% Abstract, shared by cover.tex and report.tex.
Face-forgery detectors answer ``real'' or ``fake''. A genuine photograph that has only been through a beauty app or a
retouching service is neither, and once the skin is smoothed the detectors tend to call it fake. On pixel-aligned pairs of
genuine faces and their commercially smoothed renders, sixteen of seventeen published detectors score the smoothed face as more
forged than its own original. A controlled training experiment shows that the cause is which class saw skin smoothing during
training: showing it to both classes lowers the accusation from 54 to 12 percent. We build on that recipe a face checker that
answers real, forged or beautified, marks the face part its evidence came from, and shows only sentences that code composes from
the model's output and checks before display. Restoring the marked region to its original pixels removes 94 percent of the
model's part-level forgery decisions, about twice what a post-hoc heat map achieves, and free-form descriptions written by
vision-language models fail a blind matching test, so no model is allowed to describe the image. The checker runs in a web
browser and the photograph never leaves the device.
